% ============================================================
\chapter*{Chapter 0 --- Installation and Setup}
\addcontentsline{toc}{chapter}{Chapter 0 --- Installation and Setup}
\label{ch:installation}
% ============================================================

Before you start programming, you need to set up your working environment. This chapter guides you step by step.

\section*{Installing Python}
\addcontentsline{toc}{section}{Installing Python}

\begin{intuition}
Python is a free \textbf{programming language} used by millions of scientists worldwide. Think of it as a supercharged calculator that understands instructions written in English.
\end{intuition}

\subsection*{Recommended method: Anaconda}

\textbf{Anaconda}\index{Anaconda} is a Python distribution that includes all the scientific libraries we will need.

\begin{enumerate}
\item Go to \url{https://www.anaconda.com/download}
\item Download the version for your system (Windows, macOS, or Linux).
\item Run the installer and follow the default instructions.
\item Verify the installation by opening a terminal and typing:
\end{enumerate}

\begin{codeblock}[title=\textbf{Terminal}]
\begin{minted}{bash}
python --version
\end{minted}
\end{codeblock}

\begin{codeoutput}
\begin{verbatim}
Python 3.12.4
\end{verbatim}
\end{codeoutput}

\subsection*{Lightweight alternative: Python + pip}

If you prefer a minimal installation:
\begin{enumerate}
\item Download Python from \url{https://www.python.org/downloads/}
\item Check ``Add Python to PATH'' during installation.
\item Install the required libraries:
\end{enumerate}

\begin{codeblock}[title=\textbf{Terminal}]
\begin{minted}{bash}
pip install numpy matplotlib pandas scipy seaborn scikit-learn jupyter
\end{minted}
\end{codeblock}

\section*{Jupyter Notebook}
\addcontentsline{toc}{section}{Jupyter Notebook}

\textbf{Jupyter Notebook}\index{Jupyter} is an interactive environment where you write code, see results, and add notes --- like a digital lab notebook.

\begin{codeblock}[title=\textbf{Terminal --- Launch Jupyter}]
\begin{minted}{bash}
jupyter notebook
\end{minted}
\end{codeblock}

Your browser opens automatically. Click \textbf{New $\to$ Python 3} to create a new notebook.

\begin{center}
\begin{tikzpicture}[
  box/.style={rectangle, draw, rounded corners, fill=#1, minimum width=4cm, minimum height=1cm, align=center, font=\small}
]
\node[box=blue!10] (cell1) at (0,2) {\texttt{In [1]: 2 + 3}};
\node[box=green!10] (out1) at (0,0.8) {\texttt{Out[1]: 5}};
\node[box=blue!10] (cell2) at (0,-0.5) {\texttt{In [2]: print("Hello")}};
\node[box=green!10] (out2) at (0,-1.7) {\texttt{Hello}};
\draw[-{Stealth},thick] (cell1.south) -- (out1.north);
\draw[-{Stealth},thick] (cell2.south) -- (out2.north);
\node[font=\footnotesize\itshape, right] at (3,2) {Code cell};
\node[font=\footnotesize\itshape, right] at (3,0.8) {Result};
\end{tikzpicture}
\end{center}

\begin{bestpractice}
Use \textbf{Shift + Enter} to execute a cell and move to the next one. Use \textbf{Ctrl + Enter} to execute without changing cells.
\end{bestpractice}

\section*{Text Editors}
\addcontentsline{toc}{section}{Text Editors}

To write Python scripts (\texttt{.py} files), you can use:
\begin{itemize}
\item \textbf{VS Code}\index{VS Code} (recommended): free, with Python extension.
\item \textbf{Spyder}: included in Anaconda, designed for scientists.
\item \textbf{PyCharm}: free Community edition.
\end{itemize}

\section*{Verifying the Installation}
\addcontentsline{toc}{section}{Verifying the Installation}

\begin{codeblock}[title=\textbf{Python --- First test}]
\begin{minted}{python}
import numpy as np
import matplotlib.pyplot as plt

print("Installation successful!")
print(f"NumPy version: {np.__version__}")

# Quick test: plot sin(x)
x = np.linspace(0, 2 * np.pi, 100)
plt.plot(x, np.sin(x))
plt.title("sin(x) — Your first plot!")
plt.xlabel("x")
plt.ylabel("sin(x)")
plt.grid(True)
plt.savefig('test_installation.pdf')
plt.show()
print("Plot saved!")
\end{minted}
\end{codeblock}

\begin{codeoutput}
\begin{verbatim}
Installation successful!
NumPy version: 1.26.4
Plot saved!
\end{verbatim}
\end{codeoutput}

If everything works, you are ready to start Chapter~1!

\section*{Installing R (for Chapter 9)}
\addcontentsline{toc}{section}{Installing R}

\begin{enumerate}
\item Download R from \url{https://cran.r-project.org/}
\item Install \textbf{RStudio Desktop} from \url{https://posit.co/download/rstudio-desktop/}
\item Verify by opening RStudio and typing:
\end{enumerate}

\begin{codeblock}[title=\textbf{R Console}]
\begin{minted}{r}
R.version.string
\end{minted}
\end{codeblock}

\begin{codeoutput}
\begin{verbatim}
[1] "R version 4.4.1 (2024-06-14)"
\end{verbatim}
\end{codeoutput}

\begin{errmark}
\begin{itemize}
\item \textbf{``Python is not recognized''}: you forgot to add Python to PATH. Reinstall with the option checked.
\item \textbf{``ModuleNotFoundError''}: the library is not installed. Type \texttt{pip install module\_name}.
\item \textbf{Permission denied}: on macOS/Linux, try \texttt{pip install --user module\_name}.
\end{itemize}
\end{errmark}
